% Cote 119: offprints, reports, courses, CV and letters (typed and printed part): archive cote 119, site page 273 (PDF page 274).
% typescript. Draft compared: none.
% AI-transcribed by Claude Opus 5.5 (claude-opus-5-5, adaptive thinking, effort high), Message Batches API, batch r119_full (msgbatch_01DwZNm5Y7YsKBWyAN4sFs8w), 2026-11-13; not yet checked by eye.
\scanpage{119/273}
Au cours de la 3e année, le candidat se voit confier la \unsure{responsabilite|responsabilité}
d'une tranche de cours bien délimitée. Si l'épreuve esr concluante (avis unanime
du chef de service et des étudiants) le candidat doit, l'année suivante, sous le
contrôle d'un conseiller pédagogique nommé par le département, prendre la responsa-%
bilité d'un enseignement et rédiger au moins un chapitre du cours qu'il traite.

A la fin de cette année se réunit un jury composé du conseiller pédago-%
gique, d'un rapporteur désigné par le département et d'un délégué élu des étudiants.
La plus grande attention sera donnée à des initiatives pédagogiques prises par le
candidat en accord avec son conseiller pédagogique.

Le titre de docteur ès-Sciences Mathématiques (enseignement) est décerné
sur avis unanime des membres du jury. La decision du jury est sans appel.

La prochaine réunion de la commission qui étudie cette question aura
lieu le mardi 11 juin à 9 heures en salle des délibérations. Tous les collègues
qui ont des suggestions, des critiques ou des réserves à émettre seront les
bienvenus.

-:-:-:-:-:-:-:-:-:-:-

\margin{(encircled) assistants \unsure{pas|par} à l'ancienneté} \margin{(encircled) \unsure{un peu de recherche}}

\margin{option \unsure{momentanée}}

\margin{nominations par étudiants}

\margin{thèse --} \margin{(encircled) \unsure{comp.|sup.} CNRS ?}

\margin{enseignants \unsure{parallèles} (diversification)}

\margin{(line drawn under the above)}

\margin{\unsure{recherche enseignement}} \margin{(encircled) \unsure{décision des étudiants}}

\margin{\unsure{mandarinat}}

\margin{(line drawn under the above)}

\mspage{\hand{273}}
